\documentclass[11pt,a4paper]{article}
\usepackage{fontspec}
\setmainfont{Noto Serif}
\setsansfont{Noto Sans}
\usepackage{babel}
\babelprovide[main,import]{russian}
\usepackage{amsmath,amssymb,mathtools,booktabs,geometry,fancyhdr,microtype,xcolor,hyperref}
\geometry{left=21mm,right=21mm,top=21mm,bottom=21mm}
\definecolor{ink}{HTML}{173D57}
\hypersetup{colorlinks=true,urlcolor=ink,linkcolor=ink,pdftitle={BGS-F. Мультипликативно независимые собственные фазы}}
\pagestyle{fancy}\fancyhf{}\lhead{\small BGS-F}\rhead{\small Мультипликативно независимые собственные фазы}\cfoot{\thepage}
\setlength{\headheight}{15pt}\setlength{\parindent}{1em}\setlength{\parskip}{2.5pt}
\setlength{\emergencystretch}{3em}\tolerance=3000
\newcommand{\Tr}{\operatorname{Tr}}
\newcommand{\EE}{\mathbb E}
\newcommand{\R}{\mathbb R}
\newcommand{\Z}{\mathbb Z}
\newcommand{\C}{\mathbb C}
\newcommand{\si}{\mathrm{Sine}_2}
\newcommand{\eps}{\varepsilon}
\newcommand{\sinc}{\operatorname{sinc}}
\newcommand{\sect}[1]{\section*{\sffamily\color{ink}#1}}
\newcommand{\subt}[1]{\par\medskip\noindent\textbf{#1.}\quad}
\newcommand{\T}{\mathbb T}
\newcommand{\dist}{\operatorname{dist}}
\newcommand{\Ad}{\operatorname{Ad}}
\newcommand{\zp}{\mathbb Z_7}
\newcommand{\QQ}{\mathbb Q}
\newcommand{\kk}{\overline{\mathbb Q}}
\newcommand{\ord}{\operatorname{ord}}
\newcommand{\diag}{\operatorname{diag}}
\begin{document}
\noindent{\small\sffamily\color{ink} МАТЕМАТИЧЕСКАЯ ФИЗИКА\hfill 22 НОЯБРЯ 2026}
\begin{center}{\LARGE\sffamily Мультипликативно независимые\\[3pt]собственные фазы}\\[6pt]
{\large Точный спектр при фиксированном кике на \(N=7^k\)}
\end{center}

Пусть \(U_N(t)=e^{itA_N}M_N\) с матрицами из постановки BGS-F. Исходной силе \(10^{-80}\) соответствует \(t=N/10^{80}\). Для \(N=7^k\), \(k\ge4\), положим \(R_N=180\cdot7^{k-4}\).

\subt{Линейная нижняя граница ранга}
Для любого ненулевого вещественного алгебраического \(t\) среди собственных чисел полного \(U_N(t)\) существуют \(R_N\) простых значений \(\lambda_1,\ldots,\lambda_{R_N}\), для которых
\[
      \prod_{j=1}^{R_N}\lambda_j^{a_j}\notin\overline{\mathbb Q}
         \qquad\bigl(\mathbf a\in\Z^{R_N}\setminus\{0\}\bigr).
\]
Таким образом, доказанная доля мультипликативно независимых значений по модулю алгебраических чисел равна по меньшей мере \(180/2401\). В частности, полный спектр имеет не меньше \(R_N(R_N-1)+1\) различных отношений, а произведения любого фиксированного порядка из выбранных фаз попарно различны по модулю алгебраических чисел, с точностью до перестановки сомножителей.

\subt{Выбор фаз и доказательство}
При \(t=0\) используем линейные двойные узлы \(\ell\equiv3\bmod7\). Обособленный циклотомический вклад метки \((0,7)\) фиксирует сумму коэффициентов возможной зависимости на каждом узле \(\ell\). Оставшиеся вклады меток \((0,1),(0,8),(0,-6)\) позволяют исключить шесть классов последовательно. Получаем пять независимых производных углов для каждого выбранного \(\ell\). Применение теоремы Линдемана--Вейерштрасса переносит отсутствие точных мономиальных соотношений на исходный параметр.

\subt{Следствия для усреднения по времени}
При фиксированном \(N\) вектор \((\lambda_1^m,\ldots,\lambda_{R_N}^m)\) при равномерном усреднении по \(m=1,\ldots,T\) стремится к мере Хаара на торе при \(T\to\infty\). Затем при \(k\to\infty\) выбранная часть обычного следа после деления на \(\sqrt{R_N}\) имеет стандартный комплексный гауссов закон. Для полного следа получена следующая верхняя граница плотности почти скалярных возвратов:
\[
 \limsup_{T\to\infty}\frac1T\#\left\{1\le m\le T:
        |\Tr U_N(t)^m|\ge N\left(1-\frac{90}{2401}\right)\right\}
                   \le 4e^{-R_N/16}.
\]
Это предел сначала по времени при фиксированной размерности. Предел Sine\(_2\) и парный тест при \(m\asymp N\) остаются открытыми.

\newpage
\sect{1. Перенос отношений с заданного алгебраического параметра}
Пусть \(t_*\) --- ненулевое вещественное алгебраическое число. Все элементы \(A_N,M_N\) алгебраичны, и это также верно для собственных чисел \(A_N\) и их спектральных проекторов. Поэтому коэффициенты характеристического многочлена \(U_N(t)\) принадлежат кольцу конечных сумм \(\sum_\nu a_\nu e^{itb_\nu}\), где \(a_\nu,b_\nu\) алгебраичны и \(b_\nu\) вещественны. Это кольцо сохраняется при умножении. После группировки по одинаковым \(b_\nu\) теорема Линдемана--Вейерштрасса говорит, что ненулевая сумма остаётся ненулевой в \(t_*\). В самом деле, показатели \(it_*b_\nu\) алгебраичны и различны.

При \(t\) около нуля занумеруем собственные числа аналитическими ветвями \(z_1(t),\ldots,z_N(t)\), считая кратности. Их локальное существование следует из теоремы Реллиха для эрмитова логарифма с разрезом вне спектра. Ветви, выбранные на вещественном интервале, продолжаются как аналитические функции около нуля. Обозначим через \(I_*\) идеал алгебраических полиномиальных отношений между собственными числами \(U_N(t_*)\) в некоторой первоначальной нумерации.

\subt{Одновременная нумерация}
Существует перестановка аналитических ветвей, для которой каждое отношение из \(I_*\) обращается в нуль тождественно в окрестности нуля. Иначе для каждой перестановки \(\pi\) найдётся \(P_\pi\in I_*\) с ненулевым аналитическим значением \(P_\pi(z_{\pi(1)}(t),\ldots,z_{\pi(N)}(t))\). В линейной оболочке конечного набора \(P_\pi\) над \(\QQ\) условия нулевого аналитического значения задают собственные подпространства. Их конечное объединение не заполняет оболочку. Поэтому можно выбрать один \(P\in I_*\), ненулевой как аналитическая функция при всех перестановках.

Произведение \(\prod_\pi P(X_{\pi(1)},\ldots,X_{\pi(N)})\) симметрично. Его значение на собственных числах выражается многочленом от коэффициентов характеристического многочлена. Таким образом, получаем ненулевой экспоненциальный многочлен, который всё же обнуляется в \(t_*\). Это противоречит инъективности, доказанной выше. Предыдущий перенос на линейные фазы в части о жёсткости использовал их значения при \(t=0\). Здесь сохраняются аналитические ветви целиком.

\subt{Из производных в мультипликативный ранг}
Пусть у выбранных ветвей \(z_j'(0)/(iz_j(0))=d_j\) и \(d_j\) линейно независимы над \(\QQ\). Если соответствующее произведение при \(t_*\) алгебраично, после умножения на произведение с положительными степенями оно даёт полиномиальное отношение из \(I_*\). Оно сохраняется аналитически на ветвях. Логарифмическая производная в нуле даёт \(\sum_j a_j d_j=0\), откуда \(a_j=0\) для всех \(j\). Равенство выбранного простого аналитического корня и любого другого собственного числа при \(t_*\) также было бы отношением. Поэтому ветвь, простая при всех достаточно малых ненулевых \(t\), остаётся простой в этой нумерации при \(t_*\).

\newpage
\sect{2. Стационарные вклады в производные углов}
Положим \(N=7^k\), \(m=3\cdot7^{k-1}\). Выберем \(\rho^2=2\), \(\rho\equiv3\bmod7\), и \(\gamma=3+2\rho\). В расщеплённых координатах \((a,b)=(v_1+\rho v_2,(v_1-\rho v_2)/(-2\rho))\) оператор \(M_N\) действует как дилатация с множителем \(\gamma^{-1}\) с точностью до общей фазы. На единицах есть две орбиты длины \(m\), заданные квадратностью остатка по модулю 7. На них собственные векторы определяются характером
\[
         \chi_{\ell,h}(x)=\omega(x)^h\zeta_N^{\ell\mathcal L(x)},
\quad \mathcal L(x)=\tfrac16\log(x^6),\quad \omega(3)=\zeta_3 .
\]
Здесь \(\ell\bmod7^{k-1}\), \(h\bmod3\), \(\zeta_N=e^{2\pi i/N}\). При \(\ell\equiv3\bmod7\) собственная фаза имеет кратность два, поскольку на меньших орбитах характер такого проводника отсутствует. Пара фазовых производных есть спектр сжатия \(A_N\) на это двумерное пространство.

Для данного \(\ell\) в квадратной орбите единственные стационарные метки примитивного слоя --- три противоположные пары
\[
          \pm v_1,\ \pm v_2,\ \pm v_3,\qquad
         v_1=(0,1),\quad v_2=(0,8),\quad v_3=(0,-6).
\]
Непримитивная пара есть \(\pm v_4\), где \(v_4=(0,7)\). Коэффициенты исходного кика
\[
       (c_1,c_2,c_3,c_4)
               =\frac1{4000}(200-20i,-100+10i,-100+10i,1800).
\]
Эта таблица получается подстановкой исходных меток в сравнение \(b x(x+a)=\ell a\) по модулю 7. Для \(v_4\) сначала снимается общий множитель 7. Ни одна стационарная метка в выбранном классе не соединяет разные квадратные орбиты. Нестационарные суммы обнуляются точно по выделению последнего разряда. Следовательно, сжатие диагонально в указанном орбитальном базисе.

Положим \(\Phi_i(\ell)=b(x_i+a/2)+\ell\mathcal L((x_i+a)/x_i)\), где \(x_i\) --- простой поднятый корень стационарного уравнения. Для \(i=4\) уравнение сначала делится на 7, после чего корень простой. Стационарное разложение по модулю \(7^k\) даёт
\[
 d_{\pm,\ell,h}=\frac1m\sum_{i=1}^4
   \left(c_i^\pm G_i\zeta_3^{h\mathbf1_{i<4}}\zeta_N^{\Phi_i(\ell)}
          +\overline{c_i^\pm G_i\zeta_3^{h\mathbf1_{i<4}}
                                     \zeta_N^{\Phi_i(\ell)}}\right),
\quad c_i^+=c_i,\quad c_i^-=\bar c_i .
\]
Множитель \(G_i\) ненулевой, лежит в \(\QQ(\zeta_7)\) и не зависит от \(\ell,h\) в выбранном классе. Это рациональная степень 7, при нечётном приведённом проводнике умноженная на квадратичную сумму Гаусса порядка 7. Простота корня позволяет заменить высшие члены фазы на квадрат локальной переменной. При отрицательной метке вклад комплексно сопрягается. Отражение \(x\mapsto-x\) меняет орбиты и заменяет \(c_i\) на \(\bar c_i\). Оно сохраняет конечный характер. Отношение \((x_i+a)/x_i\bmod7\) равно 4 для \(i<4\) и 1 для \(i=4\). Поэтому указанная степень \(\zeta_3\) учтена точно.

\newpage
\sect{3. Циклотомические координаты и подъёмы}
Пусть \(F=\QQ(i,\zeta_3)\), \(K=F(\zeta_7)\), \(D=7^{k-1}\) и \(B=7^{k-2}\). Базис \(K(\zeta_N)\) над \(K\) образован \(1,\zeta_N,\ldots,\zeta_N^{D-1}\). Это следует по степени из взаимной простоты 12 и \(7^k\). Поэтому независимо выделяются координаты с заданными остатками показателя по любому делителю \(D\). В частности, остатки \(\Phi_i\bmod49\) для \(i=1,2,3,4\) равны \(10,3,17,14\) соответственно. Они и четыре противоположных остатка различны. Так можно выделить вклад любой положительной метки без её сопряжённой части.

Для каждого выбранного представителя \(r\bmod B\), \(r\equiv3\bmod7\), используем шесть подъёмов \(\ell=r+uB\), \(0\le u\le5\). В каждой семье фаз
\[
  \Phi_i(r+uB)\equiv\Phi_i(r)
                  +uB\mathcal L((x_i(r)+a)/x_i(r))\pmod{7^k}.
\]
Действительно, \(x_i(r+uB)-x_i(r)\) кратно \(B\) по простому уравнению Гензеля. Разложение фазы около стационарной точки исключает линейный член по \(x\). Остаток квадратичен по \(B\) и исчезает по модулю \(7^k\), поскольку \(k\ge4\). Член по \(\ell\) равен записанному логарифму. Этот логарифм кратен 7. Если его частное на 7 есть единица, шесть подъёмов отличаются по множителям \(\zeta_7^{uv_i}\), где \(v_i\ne0\bmod7\), и имеют одну координату по модулю \(D\). Эти шесть степеней линейно независимы над \(F\), поскольку единственная зависимость между всеми семью степенями над \(F\) пропорциональна сумме семи корней единицы.

\subt{Различимость представителей в одном классе}
Если \(r,r'\) в одном классе по модулю 49 и \(\mathcal L((x_i(r)+a)/x_i(r))/7\) обратим, то
\[
              \nu_7\bigl(\Phi_i(r')-\Phi_i(r)\bigr)
                         =\nu_7(r'-r)+1 .
\]
Линейный член по \(r'-r\) имеет такой порядок. Квадратичная ошибка имеет порядок не меньше \(2\nu_7(r'-r)\), который строго больше. Для непримитивной метки \(v_4\) верно то же для любых \(r,r'\equiv3\bmod7\). Здесь вся фаза имеет множитель 7, квадратичная ошибка имеет порядок не меньше \(1+2\nu_7(r'-r)\), а обратимость следует из
\[
         \mathcal L((x_4+a)/x_4)/7\equiv(a/7)/x_4\equiv5\pmod7.
\]
Таким образом, выбранные различные представители \(r\bmod B\) дают различные координаты по модулю \(D\) в каждом нужном обратимом классе. Остатки \(\Phi_i\bmod49\) стабильны при изменении \(\ell\) на кратное 7, а остатки по модулю 343 стабильны при изменении на кратное 49. Обе стабильности следуют из указанного разложения и делимости первой производной на 7. Остаток частного \(\Phi_i^{\prime}/7\) вычисляется через \(x_i\bmod49\).

\newpage
\sect{4. Пять независимых углов на каждом выбранном узле}
Берём представителей \(0\le r<B\) в шести классах \(r\equiv3,10,17,24,38,45\bmod49\), с каждым из шести подъёмов. Из шести производных на фиксированном \(\ell\), двух ветвей при каждом из трёх \(h\), убираем ветвь \((- ,2)\). Пусть \(n_{\ell,h,\pm}\in\QQ\) --- коэффициенты возможной зависимости. Отсутствующий коэффициент считается нулевым.

Сначала выделяем циклотомические координаты \(\Phi_4\equiv14\bmod49\). Для разных представителей \(r\) они различны по модулю \(D\). При каждом представителе шесть подъёмов разделяются независимостью степеней \(\zeta_7\) над \(F\). Непримитивный коэффициент и характер совпадают для всех \(h,\pm\). Значит, для каждого \(\ell\) выполняется
\[
                       \sum_{h,\pm}n_{\ell,h,\pm}=0 .                 \tag{1}
\]

Для примитивного вклада используем таблицу пар \((\Phi_i\bmod343,\Phi_i'/7\bmod7)\):
\begin{center}\small
\begin{tabular}{crrr}
\toprule
\(r\bmod49\)&\(i=1\)&\(i=2\)&\(i=3\)\\
\midrule
3  &\((206,5)\)&\((150,2)\)&\((66,1)\)\\
10 &\((206,2)\)&\((3,6)\)&\((213,5)\)\\
17 &\((59,6)\)&\((52,3)\)&\((213,2)\)\\
24 &\((108,3)\)&\((297,0)\)&\((66,6)\)\\
31 &\((10,0)\)&\((52,4)\)&\((115,3)\)\\
38 &\((108,4)\)&\((3,1)\)&\((17,0)\)\\
45 &\((59,1)\)&\((150,5)\)&\((115,4)\)\\
\bottomrule
\end{tabular}
\end{center}
В таблице выделяем классы в следующем порядке. Первая координата задаёт класс, вторая номер вклада.
\[
   (17,2),\ (45,3),\ (10,3),\ (3,1),\ (38,2),\ (24,1).
\]
 В каждый момент нужный фазовый остаток по модулю 343 уникален среди оставшихся классов для данной метки. Разные метки уже разделены по модулю 49. Частное производной на 7 обратимо в каждой используемой строке. Поэтому на каждом шаге отдельно выделяем каждый представитель и каждый подъём точно тем же аргументом с шестью степенями \(\zeta_7\). После снятия ненулевого общего множителя получаем
\[
 \sum_{h=0}^2\zeta_3^h
       (c_i n_{\ell,h,+}+\bar c_i n_{\ell,h,-})=0 .                    \tag{2}
\]
Над \(\QQ(i)\) единственное направление зависимости трёх степеней \(\zeta_3\) имеет равные коэффициенты. Так как \(c_i,\bar c_i\) независимы над \(\QQ\), коэффициенты при каждой ветви постоянны по \(h\). Один из них равен нулю по выбору, второй равен нулю по (1). Это удаляет весь рассматриваемый класс и даёт индукцию.

Число выбранных производных равно \(6\cdot7^{k-4}\cdot6\cdot5=R_N\). Все выбраны из ветвей, простых при малом ненулевом \(t\). Для ветви \((+,2)\) различие с исключённой \((-,2)\) тоже следует из ненулевой компоненты \((c_1-\bar c_1)G_1\) в остатке 10 по модулю 49. Применяем перенос из раздела 1 к этому набору.

\newpage
\sect{5. Длинное время и счёт точных отношений}
Пусть \(\lambda_1,\ldots,\lambda_R\) --- найденные собственные числа точного оператора в алгебраическом параметре \(t\), \(R=R_N\). Для каждого \(\mathbf a\in\Z^R\setminus\{0\}\) значение \(\prod_j\lambda_j^{a_j}\) трансцендентно. Поэтому средней геометрической прогрессии с этим знаменателем соответствует предел ноль. По критерию Вейля вектор \((\lambda_1^m,\ldots,\lambda_R^m)\) при равномерном выборе \(m\in\{1,\ldots,T\}\) сходится к произведению \(R\) равномерных мер на окружности. Та же сходимость верна при любом фиксированном ненулевом шаге арифметической прогрессии, поскольку никакая ненулевая степень соответствующего произведения не равна единице. Для фиксированных \(N\) и шага сходимость на непрерывных функциях равномерна по началу окна: любой фиксированный ненулевой характер оценивается геометрической суммой, не зависящей по модулю от начального индекса.

\subt{След выбранного спектрального проектора}
Пусть \(P_N\) --- спектральный проектор на выбранные простые фазы. Тогда \(R^{-1/2}\Tr(P_NU_N^m)\) после первого предела \(T\to\infty\) имеет закон суммы \(R\) независимых единичных комплексных фаз, делённой на \(\sqrt R\). При втором пределе \(k\to\infty\) получается закон с плотностью \(\pi^{-1}e^{-|z|^2}\). Например, характеристическая функция в направлении \(u\in\C\) для одного слагаемого равна \(1-|u|^2/(4R)+O_u(R^{-3/2})\) по симметрии окружности. Возведение в степень \(R\) даёт нужный предел.

\subt{Почти скалярные возвраты полного оператора}
Условие \(|\Tr U_N^m|\ge N-R/2\) влечёт \(|\Tr(P_NU_N^m)|\ge R/2\) по треугольнику, поскольку оставшихся фаз \(N-R\). Для независимых равномерных фаз \(Z_j\) действительная и мнимая части центрированы и симметричны. Для каждой из них моментная функция ограничена \(e^{s^2/2}\) из \(\cosh x\le e^{x^2/2}\). При \(|\sum Z_j|\ge R/2\) какая-то координата имеет модуль не меньше \(R/(2\sqrt2)\). Экспоненциальная оценка четырёх хвостов даёт вероятность не более \(4e^{-R/16}\). Применяем верхнюю границу слабого предела к замкнутому событию. Получается заявленная верхняя плотность времён. При \(t=N/10^{80}\) доказательство относится к полной исходной матрице.

\subt{Количественная комбинаторика}
Все \(R(R-1)\) направленных отношений \(\lambda_i/\lambda_j\), \(i\ne j\), различны между собой даже по модулю \(\overline{\QQ}^{\,*}\). В противном случае возник бы ненулевой вектор показателей. Добавляя единицу, получаем не менее \(R(R-1)+1\) частот в полном спектре разностей. Для произведений ровно \(d\) сомножителей из выбранных фаз число различных классов равно \(\binom{R+d-1}{d}\). Эти утверждения касаются точного равенства. Плотность близких частот в интервале длины \(1/N\) из них не следует.

\newpage
\sect{6. Воспроизведение и область применимости}
\subt{Конечная арифметическая проверка}
Команда \texttt{python3 relation\_rank/check.py --verify} воспроизводит таблицу стационарных фаз, коэффициенты разделения по модулю 7 и порядок исключения шести классов. Дополнительно проверяются подъёмы до \(k=7\), разделимость циклотомических координат и ненулевой малый минор над \(\mathbb F_{101}\). В стационарных вычислениях используется целочисленный логарифм по модулю степени 7 и однозначное поднятие простого корня. Файл \texttt{relation\_rank/data/certificates.json} содержит точные данные.

\subt{Прямая проверка ранга}
Скрипт \texttt{relation\_rank/check\_rank.py} формирует коэффициентные векторы самих фазовых производных в тензорном базисе по \(i,\zeta_3,\zeta_{7^k}\). Нормировка устраняется общим целым множителем. Квадратичные гауссовы множители разворачиваются в явные суммы степеней \(\zeta_7\). Затем ранг вычисляется над \(\mathbb F_{101}\) без использования порядка исключения из раздела 4. Проверены значения
\begin{center}
\begin{tabular}{rrr}
\toprule
\(k\)&\(N\)&Ранг над \(\mathbb F_{101}\)\\
\midrule
4&2401&180\\
5&16807&1260\\
6&117649&8820\\
\bottomrule
\end{tabular}
\end{center}
Поскольку коэффициентные векторы целочисленные, конечнополевой ранг даёт нижнюю границу рационального ранга при этих размерностях. Для произвольного \(k\ge4\) используется циклотомическая аргументация, изложенная выше.

\subt{Граница для исходного предела}
На всех этих размерностях точные резонансы среди выбранного набора исключены при неизменной силе кика. Доказательство усреднения использует геометрические множители с произвольными ненулевыми расстояниями \(|1-\prod_j\lambda_j^{a_j}|\). Оно не даёт оценку этих знаменателей, достаточную для окна времени порядка \(N\). Поэтому результат относится к точной арифметике полного спектра и последовательному пределу по времени. Для решения BGS-F всё ещё требуется контролировать приближённые совпадения на микроскопическом масштабе.

\subt{Входные утверждения}
Дилатационная форма, базис характеров и точная стационарная фаза описаны в главах о кратностях и законе расщепления, входящих в сводный отчёт. Теорема Реллиха изложена в 
\href{https://www.mat.univie.ac.at/~armin/publ/handbook_paper.pdf}{обзоре о возмущениях многочленов и линейных операторов, теорема 2.1}. Инъективность алгебраической специализации опирается на 
\href{https://encyclopediaofmath.org/wiki/Lindemann_theorem}{теорему Линдемана--Вейерштрасса}.

\end{document}
